\chapter{Methodology}
\label{ch:methodology}
% ============================================================

\todo{Describe your study design in detail. Include ethical considerations,
participant recruitment, and data collection instruments.}

\section{Study Design Overview}

The field study follows a three-phase model:

\begin{enumerate}
  \item \textbf{Phase 1 – Baseline (without AI assistance):} Participants
    produce journalistic texts without any generative AI tools, establishing
    baseline metrics.
  \item \textbf{Phase 2 – Intervention (with AI assistance):} Participants use
    the custom web application with an integrated generative AI interface under
    comparable task conditions.
  \item \textbf{Phase 3 – Qualitative Interview:} A semi-structured interview
    captures subjective perceptions, usage strategies, and critical reflections.
\end{enumerate}

\section{Web Application}

\subsection{Architecture}
\subsection{Logging and Metrics Collection}

\section{Participants}

\section{Data Collection}

\subsection{Quantitative Metrics}
\subsection{Qualitative Instruments}

\section{Data Analysis}

\subsection{Efficiency Analysis}
\subsection{Writing Style Analysis}

\section{Hypotheses}

Based on the literature, the following hypotheses are proposed:

\begin{description}
  \item[$H_1$] AI assistance leads to a statistically significant reduction in
    task completion time (RQ1).
  \item[$H_2$] AI-assisted texts exhibit measurable stylistic homogenisation
    across participants (RQ2).
\end{description}

\section{Ethical Considerations}